\documentclass[11pt, a4paper]{article}
\usepackage[a4paper, margin=2.5cm]{geometry}
\usepackage{iftex}
\ifPDFTeX
  \usepackage[utf8]{inputenc}
  \usepackage[T1]{fontenc}
  \usepackage[english, spanish]{babel}
\else
  \usepackage{fontspec}
  \usepackage[english, spanish]{babel}
  \babelprovide[import]{spanish}
  \IfFontExistsTF{Noto Sans}{\babelfont{rm}{Noto Sans}}{}
\fi

\usepackage{amsmath}
\usepackage{booktabs}
\usepackage{tabularx}
\usepackage{enumitem}
\usepackage{xcolor}
\usepackage{mdframed}
\usepackage[hidelinks]{hyperref}

\definecolor{usbblue}{RGB}{13, 110, 253}
\definecolor{darkgray}{RGB}{50, 50, 50}
\definecolor{lightgray}{RGB}{245, 245, 245}
\definecolor{quoteframe}{RGB}{200, 200, 200}

\newmdenv[
  topline=false,
  bottomline=false,
  rightline=false,
  leftline=true,
  linecolor=quoteframe,
  linewidth=3pt,
  backgroundcolor=lightgray,
  innerleftmargin=12pt,
  innerrightmargin=12pt,
  innertopmargin=8pt,
  innerbottommargin=8pt
]{guionquote}

\pagestyle{plain}

\begin{document}

\begin{center}
    {\LARGE \textbf{Guion de Defensa — Informe de Pasantía}}\\[1.5em]
    {\Large \textbf{Título:} Diseño e implementación de un sistema IoT con Microcontroladores para el Monitoreo Remoto de Equipos Fiscales}\\[1em]
    \textbf{Autora:} Lorena Del Valle Rojas Chams El Dein — C.I. 19-10469 — Ingeniería Electrónica, Universidad Simón Bolívar\\
    \textbf{Institución:} The Factory HKA\\
    \textbf{Tutor académico:} Prof. Miguel Díaz\\
    \textbf{Tutores industriales:} Ing. Pablo Arguinzones y Ing. Luis Muñoz\\
    \textbf{Formato:} 23 diapositivas $\cdot$ $\sim$3.300 palabras habladas $\cdot$ \textbf{objetivo 30–32 min} + ronda de preguntas
\end{center}

\vspace{1em}

\begin{guionquote}
\textbf{Versión condensada (30–32 min).} La versión extensa quedó respaldada en \texttt{Speech\_Defensa\_Pasantias\_respaldo\_largo.md}. Este guion ya incorpora las correcciones del libro v2: el enlace local es \textbf{síncrono (100 kHz)}, distinto del \textbf{RS-232/UART asíncrono} inicial; el conflicto de concurrencia es un \textbf{acceso concurrente al bus compartido}, arbitrado con un \textbf{mutex de software de FreeRTOS}; el objetivo del servidor MQTT figura como \textbf{parcialmente logrado}; y el \textbf{Bloque 3 está reenfocado a resultados}.
\end{guionquote}

\section*{Cómo usar este guion}

\begin{itemize}[label=-]
    \item El texto entre \textbf{negrita y corchetes} son acotaciones para la ponente (\texttt{[PAUSA]}, \texttt{[SEÑALAR]}, \texttt{[MOSTRAR]}). \textbf{No se leen en voz alta.}
    \item El discurso va en \textbf{primera persona del plural} para toda la exposición técnica; el singular se reserva para la presentación personal, los agradecimientos y la reflexión final.
    \item Los tiempos son acumulativos, calculados a $\sim$115 palabras por minuto (tu ritmo medido). Si en el ensayo vas más rápido, tendrás margen; si te atascas, recorta la diapositiva 9.
    \item Las frases de transición cierran cada bloque: son clave para la continuidad.
\end{itemize}

\section*{Mapa de diapositivas (23 en total)}

\begin{table}[htbp]
\centering
\small
\begin{tabularx}{\textwidth}{c X c}
\toprule
\textbf{\#} & \textbf{Diapositiva} & \textbf{Bloque} \\
\midrule
1 & Portada: título, autora, tutores, institución & 1 \\
2 & Contexto: The Factory HKA y el parque de impresoras fiscales & 1 \\
3 & Planteamiento del problema & 1 \\
4 & Objetivos: general y específicos & 1 \\
5 & Arquitectura IIoT de tres niveles (Edge / Gateway / Cloud) & 2 \\
6 & Nivel Edge: captura del \textit{Status S1} por interfaz síncrona & 2 \\
7 & Nivel Gateway: ESP32, ESP-IDF y kernel concurrente (FreeRTOS) & 2 \\
8 & Nivel Cloud: MQTT v5.0 / TLS 1.2 y Django + HTMX & 2 \\
9 & Justificación tecnológica: Arduino vs. ESP-IDF $\cdot$ HTTP vs. MQTT & 2 \\
10 & Metodología de validación: banco Hardware-in-the-Loop (HIL) & 3 \\
11 & \textbf{Resultado 1} — Captura íntegra e integridad de la telemetría & 3 \\
12 & \textbf{Resultado 2} — Robustez e inmunidad del bus local & 3 \\
13 & \textbf{Resultado 3} — Eficiencia de red: ahorro del 80,2 \% & 3 \\
14 & \textbf{Resultado 4} — Resiliencia: cero pérdida de datos & 3 \\
15 & \textbf{Resultado 5} — Latencia End-to-End ($T_{E2E}$) & 3 \\
16 & \textbf{Resultado 6} — Actualización remota OTA de doble objetivo & 3 \\
17 & \textbf{Resultado 7} — Rendimiento global: memoria, energía y trazabilidad & 3 \\
18 & \textbf{Resultado 8} — Comparativa: estado previo vs. estado actual & 3 \\
19 & Verificación de cumplimiento de objetivos & 4 \\
20 & Aporte y beneficio para la organización & 4 \\
21 & Limitaciones encontradas y trabajo futuro & 4 \\
22 & Escalabilidad y recomendaciones & 4 \\
23 & Cierre, agradecimientos y paso a preguntas & 5 \\
\bottomrule
\end{tabularx}
\end{table}

\newpage

\section*{BLOQUE 1 — Introducción y Marco del Problema}
\textbf{Tiempo estimado: 3,6 min (acumulado 0:00 – 3:36) $\cdot$ Diapositivas 1 a 4}

\subsection*{Diapositiva 1 — Portada (0:00 – 0:35)}

[\textbf{Diapositiva en pantalla. Habla pausado, mirando al jurado.}]

Primero que nada, buenos días. Nos gustaría comenzar extendiéndoles un cordial saludo a todos los presentes y presentándonos: mi nombre es Lorena Rojas. El día de hoy estamos aquí para presentarles la defensa de mi Proyecto de Pasantía, titulado \textit{Diseño e implementación de un sistema IoT con Microcontroladores para el Monitoreo Remoto de Equipos Fiscales}.

Este trabajo fue desarrollado en la empresa The Factory HKA, bajo la tutoría académica del profesor Miguel Díaz y la tutoría industrial del ingeniero Pablo Arguinzones, a quienes agradezco desde ya su acompañamiento durante todo el proceso. [\textbf{PAUSA BREVE}]

\subsection*{Diapositiva 2 — Contexto (0:35 – 1:29)}

[\textbf{SEÑALAR la imagen del parque de impresoras o el logo de HKA.}]

The Factory HKA es una empresa líder en el sector de equipos fiscales en Venezuela. Para entender el problema, veamos cómo opera realmente una impresora fiscal en campo: no es un simple periférico de impresión, es una \textbf{unidad de control autorizada por el SENIAT} bajo la Providencia 0141. Cada factura que emite queda grabada en una memoria fiscal protegida —nadie puede alterarla—, y al cierre de la jornada el equipo genera el \textit{Reporte Z}: el documento con validez legal que consolida la recaudación tributaria del día.

Ese reporte no se lleva en papel: lo transmite al SENIAT un \textbf{módulo de comunicación que el propio equipo fiscal lleva incorporado} —y ese módulo lo fabrica la empresa—. De ahí nace la oportunidad de esta pasantía: el hardware que saca datos del equipo ya existe en campo; la intención del proyecto es aprovechar ese mismo módulo, no solo para rendir cuentas al fisco, sino también para servir a la empresa.

\subsection*{Diapositiva 3 — Planteamiento del problema (1:29 – 2:29)}

[\textbf{PAUSA. Mira al jurado.}]

El laboratorio de desarrollo y la unidad de soporte de The Factory HKA enfrentan una limitación concreta: hoy ese módulo solo rinde cuentas al SENIAT. El estado operativo del equipo —el registro \textit{Status S1}, los errores, los contadores— sigue siendo invisible entre un reporte y el siguiente: conocerlo exige la presencia física del técnico o una conexión serial punto a punto.

Esto genera tres consecuencias. Primero, dificulta la detección temprana de errores críticos, como \textit{memoria fiscal llena} o fallos mecánicos, que pueden interrumpir la actividad comercial de forma imprevista. Segundo, la visualización es fragmentada: los técnicos observan un equipo a la vez, sin visión global del parque. Y tercero, no existe un historial de auditoría consultable.

[\textbf{PAUSA}] En síntesis: \textbf{equipos fiscales con datos valiosos, pero aislados, sin telemetría centralizada y sin trazabilidad.}

\subsection*{Diapositiva 4 — Objetivos (2:29 – 3:36)}

[\textbf{SEÑALAR el objetivo general.}]

Frente a esto, el objetivo general fue: \textbf{diseñar, implementar y validar un sistema basado en microcontroladores ESP32 que permita el monitoreo remoto y en tiempo real de los datos operativos de equipos fiscales, garantizando la confiabilidad y la robustez del sistema de comunicación.}

[\textbf{SEÑALAR los objetivos específicos. Agrupados.}]

Los ocho objetivos específicos se agrupan en tres frentes:

\begin{itemize}[label=-]
    \item \textbf{Diseño e integración:} definir la arquitectura del sistema IoT y desarrollar el firmware embebido, integrándolo con las interfaces del equipo fiscal.
    \item \textbf{Comunicación y visualización:} implementar la mensajería MQTT, la capa de recepción y distribución de los datos, y el dashboard web en tiempo real.
    \item \textbf{Validación y documentación:} evaluar confiabilidad, latencia y robustez ante fallos, y documentar todo el proceso.
\end{itemize}

[\textbf{TRANSICIÓN al Bloque 2}] Con el problema y los objetivos claros, pasamos a mostrarles cómo diseñamos la solución.

\hrulefill

\section*{BLOQUE 2 — Arquitectura y Diseño de la Solución}
\textbf{Tiempo estimado: 7,2 min (acumulado 3:36 – 10:50) $\cdot$ Diapositivas 5 a 9}

\subsection*{Diapositiva 5 — Arquitectura IIoT de tres niveles (3:36 – 4:51)}

[\textbf{SEÑALAR el diagrama. Recórrelo de arriba hacia abajo con el puntero.}]

La solución es una arquitectura distribuida de tres niveles, alineada con los estándares industriales del Internet Industrial de las Cosas. La denominamos \textit{MT-IoT}: Módulo de Transmisión IoT.

[\textbf{SEÑALAR cada nivel.}]

En el \textbf{Nivel de Borde} está la impresora fiscal HKA80, donde se origina la información crítica: los registros \textit{Status S1} y el diagnóstico \textit{SV2}, capturados por una interfaz síncrona transparente. En el \textbf{Nivel de Pasarela} está el Módulo de Transmisión IoT, basado en el ESP32: actúa como maestro de la interfaz local y como cliente seguro hacia la red. Y en el \textbf{Nivel Cloud} están el \textit{broker} HiveMQ y el servidor Django, que gestionan la recepción de la telemetría y su visualización en tiempo real.

[\textbf{SEÑALAR las flechas.}]

El flujo de datos va de la impresora al ESP32 por el bus local; del ESP32 a la nube por Wi-Fi con TLS; y de la nube al dashboard por WebSockets seguros.

\subsection*{Diapositiva 6 — Nivel Edge: captura del \textit{Status S1} (4:51 – 6:31)}

[\textbf{SEÑALAR el detalle del bus síncrono.}]

El nivel de borde fue el más delicado, porque es donde se toca la información fiscal directamente.

Conviene aclarar una distinción importante. La impresora HKA80 se comunica con parámetros obligatorios del protocolo —9600 baudios, 8 bits de datos, paridad par y 1 bit de parada— sobre el estándar asíncrono RS-232. Ese fue el punto de partida. Sin embargo, el enlace definitivo de la solución \textbf{no es asíncrono, sino síncrono}: migramos la interconexión física local hacia una \textbf{interfaz serie síncrona dedicada operando a una frecuencia de reloj de 100 kHz}.

[\textbf{SEÑALAR el diagrama maestro-esclavo.}]

La diferencia es de fondo: el estándar asíncrono no comparte señal de reloj y cada extremo temporiza con su propio reloj, así que cualquier diferencia entre ambos se acumula a lo largo de la trama —esa acumulación es la deriva—. En la interfaz síncrona, el ESP32 genera el reloj como maestro y el periférico muestrea al ritmo de ese mismo reloj; al compartir una única referencia de tiempo, no hay desajuste posible y la transmisión es determinista.

[\textbf{SEÑALAR el detalle de la trama.}]

Para blindar las tramas tributarias validamos con redundancia longitudinal (LRC): operaciones XOR acumulativas contrastadas contra el \textit{checksum} del periférico. Si la trama es válida, respondemos con un \textit{ACK}; si no, con un \textit{NAK}, reiniciando el bus para impedir que datos corruptos se propaguen.

\subsection*{Diapositiva 7 — Nivel Gateway: ESP32, ESP-IDF y kernel concurrente (6:31 – 8:07)}

[\textbf{SEÑALAR el diagrama de los dos núcleos.}]

Para el nodo de transmisión seleccionamos la tarjeta ESP32 DevKitC V4, de doble núcleo: mientras un núcleo sondea la impresora cada 5 segundos, el segundo administra la pila de red y la mensajería MQTT. Así, las latencias de internet no bloquean la captura local de datos.

[\textbf{SEÑALAR el diagrama de exclusión mutua.}]

El corazón del diseño es el \textbf{kernel concurrente}. Sobre FreeRTOS, la tarea de adquisición fiscal y la tarea de red se ejecutan de forma independiente, y la coordinación entre ambas se resuelve con \textbf{exclusión mutua}: dos semáforos \textit{mutex} de software que custodian los recursos compartidos del sistema.

\begin{itemize}[label=-]
    \item El primer \textit{mutex}, \texttt{mqtt\_data\_mutex}, protege la \textbf{estructura de telemetría compartida}. Garantiza que la tarea de red lea siempre un conjunto de campos coherente y completo, y nunca una lectura a medio actualizar.
    \item El segundo, \texttt{bus\_local\_mutex}, arbitra el \textbf{acceso al bus físico}. Asegura que una sola tarea energice el canal de comunicación a la vez —ya sea el sondeo periódico o una actualización OTA—, preservando la integridad de cada transacción.
\end{itemize}

\subsection*{Diapositiva 8 — Nivel Cloud: MQTT v5.0 / TLS y Django + HTMX (8:07 – 9:46)}

[\textbf{SEÑALAR el esquema de tópicos MQTT.}]

En la capa de red adoptamos el estándar \textbf{MQTT v5.0 sobre TLS 1.2}, hacia la infraestructura de HiveMQ Cloud. Se optó por MQTT por su mínima sobrecarga de cabecera y su modelo publicación-suscripción, y por la versión 5.0 por dos capacidades clave:

\begin{itemize}[label=-]
    \item Las \textbf{\textit{User Properties}}: metadatos en la cabecera del protocolo —modelo, versión de firmware, serial y RSSI— que el servidor lee sin deserializar el cuerpo del mensaje.
    \item El \textbf{\textit{Last Will and Testament}}: una "última voluntad" que el dispositivo registra al conectar; si pierde la conexión abruptamente, el \textit{broker} la publica y el sistema sabe que el equipo cayó.
\end{itemize}

La seguridad usa el puerto 8883 sobre SSL y un \textit{bundle} de certificados que valida la cadena de confianza.

[\textbf{SEÑALAR el diagrama del servidor.}]

En la nube construimos un ecosistema web con Django y HTMX. La pieza clave es un \textit{worker} MQTT independiente que corre en segundo plano y descompone cada trama en eventos atómicos: cambios de estado, cortes Z, facturas, reconexiones y actualizaciones. Para la interfaz integramos el cliente MQTT de JavaScript directamente en el navegador, sobre el canal seguro del puerto 8884, y con HTMX logramos una experiencia de página única.

\subsection*{Diapositiva 9 — Justificación tecnológica (9:46 – 10:50)}

[\textbf{SEÑALAR la tabla comparativa. Sé breve.}]

Permítannos justificar las dos decisiones tecnológicas más importantes.

\textbf{Arduino frente a ESP-IDF.} El modelo de ejecución de Arduino es un bucle secuencial mono-hilo; un \texttt{delay()} bloqueante en la pila Wi-Fi arriesgaba sobrepasar la ventana temporal del firmware legal del equipo. Por eso optamos por el framework nativo \textbf{ESP-IDF v5.5.2}, que nos dio gestión explícita de memoria en C y asignación determinista de tareas con FreeRTOS.

\textbf{HTTP frente a MQTT.} Una arquitectura HTTP REST tradicional tiene una sobrecarga de cabecera enorme para un mensaje que cambia cada 5 segundos. MQTT traslada los metadatos a la cabecera y aligera el cuerpo; más adelante les mostraremos la medición cuantitativa de ese ahorro.

[\textbf{TRANSICIÓN al Bloque 3}] Y con esto llegamos al bloque central: la validación y los resultados experimentales.

\hrulefill

\section*{BLOQUE 3 — Validación y Resultados Experimentales}
\textbf{Tiempo estimado: 11,9 min (acumulado 10:50 – 22:42) $\cdot$ Diapositivas 10 a 18}

\begin{guionquote}
\textbf{Nota de enfoque:} este bloque presenta \textbf{qué se obtuvo y cómo se comprobó}, no los problemas del camino. Esos hallazgos ya están en el informe; aquí solo se mencionan si son indispensables para entender un resultado.
\end{guionquote}

\subsection*{Diapositiva 10 — Metodología de validación: banco HIL (10:50 – 12:17)}

[\textbf{SEÑALAR la foto del banco de pruebas.}]

Antes de los resultados, explicamos brevemente \textbf{cómo se validó el sistema}.

Un equipo fiscal \textbf{no admite ensayos destructivos}: inducir cortes de energía o corromper tramas sobre una impresora real comprometería registros tributarios. Y había además una restricción práctica: el firmware del periférico seguía en desarrollo por parte de la empresa.

Para conciliar este tema, se construyó un entorno \textbf{Hardware-in-the-Loop} con dos niveles de emulación. El \textbf{primero} es un emulador en software, \texttt{impresora.py}, desarrollado con \texttt{pyserial} y los mismos parámetros eléctricos del enlace fiscal: implementa la máquina de estados de recepción, calcula el LRC y responde con las tramas S1, SV2 y los acuses \textit{ACK} y \textit{NAK}. El \textbf{segundo} eleva la fidelidad: un ESP32 esclavo que emula íntegramente al periférico sobre el mismo bus y replica el protocolo completo de reescritura remota.

Contra ese banco validamos la OTA de extremo a extremo e inyectamos fallos de forma determinista —paquetes fuera de secuencia, \textit{checksums} alterados, cortes de enlace— sin comprometer un equipo real. Y hoy cerramos el ciclo: con la \textbf{impresora fiscal HKA80} ya disponible, la demostración de esta defensa se hace sobre el equipo real junto al ESP32.

[\textbf{PAUSA}] Todos los resultados que verán fueron medidos sobre este banco de pruebas.

\subsection*{Diapositiva 11 — Resultado 1: captura íntegra e integridad de la telemetría (12:17 – 13:51)}

[\textbf{SEÑALAR la captura de la consola serial.}]

El primer resultado es la \textbf{captura íntegra y transparente del registro \textit{Status S1}}. Como ven en la traza, el módulo aísla los registros de la Memoria de Trabajo del equipo fiscal: el RIF del contribuyente, el número de registro de la máquina y los acumuladores de ventas diarias, y confirma el estado mecánico y lógico —modo fiscal y en espera, estado \texttt{0x60}, sin errores—.

Y esa captura es \textbf{transparente}: el sondeo periódico \textbf{no altera la máquina de estados del periférico ni interrumpe las transacciones en curso} —el módulo observa, no interviene—. El análisis de cada trama, además, corre sobre memoria estática, sin reservas dinámicas en la ruta crítica, con un tiempo de proceso acotado y constante.

En la misma línea, el segundo resultado: \textbf{la telemetría que llega a la nube es siempre válida y deserializable}. El campo de versión del periférico se entrega íntegro —\texttt{020507GD00}— y el \textit{worker} de Django sanea los caracteres de control antes de persistir, de modo que ningún residuo llega a la base de datos. Aquí vale mencionar un detalle del protocolo que detectamos en el banco: la trama SV2 incluye un salto de línea invisible justo antes del byte de fin de trama, que desplazaba la extracción un byte —un error \textit{off-by-one}— y corrompía el \textit{payload}. Lo resolvimos delimitando dinámicamente el fin de la cadena e inyectando el terminador nulo de forma explícita; por eso hoy el campo llega limpio.

\subsection*{Diapositiva 12 — Resultado 2: robustez e inmunidad del bus local (13:51 – 15:06)}

[\textbf{SEÑALAR la tabla de los tres niveles.}]

El segundo resultado es la \textbf{robustez del bus local}, el eslabón físicamente más expuesto del sistema. La obtuvimos con una defensa escalonada en tres niveles:

\begin{itemize}[label=-]
    \item \textbf{Defensa temporal:} cada operación se ejecuta bajo ventanas estrictas —300 milisegundos para el \textit{ACK} de reescritura, 2000 para las respuestas de estado y 3000 para el reporte S1— con una regla de máximo 3 reintentos.
    \item \textbf{Sanitización e integridad:} limpieza determinista de los búferes antes de cada transacción y validación LRC, que descarta cualquier trama alterada y responde con \textit{NAK}.
    \item \textbf{Reconstrucción estructural:} si los errores superan el umbral, el firmware \textbf{desmonta y reancla por completo el bus local}, sin reiniciar el microcontrolador.
\end{itemize}

[\textbf{SEÑALAR el resultado final.}]

El resultado medido: ante la desconexión física del periférico o la inyección de tramas corruptas, el sistema \textbf{se recuperó automáticamente}, sin reiniciar el SoC y \textbf{sin que una sola trama inválida alcanzara la base de datos}.

\subsection*{Diapositiva 13 — Resultado 3: eficiencia de red, ahorro del 80,2 \% (15:06 – 15:57)}

[\textbf{SEÑALAR la tabla comparativa. Dila con calma.}]

El tercer resultado es la \textbf{eficiencia de red}. Medimos el consumo por cada ciclo de telemetría frente a una arquitectura HTTP tradicional.

[\textbf{SEÑALAR fila por fila.}]

\begin{itemize}[label=-]
    \item \textbf{Cabecera:} de unos 450 bytes en HTTP a 12 bytes en MQTT. Una reducción del \textbf{97,33 \%}.
    \item \textbf{Cuerpo del mensaje:} de 328 a 142 bytes. Un \textbf{56,70 \%} menos.
    \item \textbf{Consumo total por trama:} de \textbf{778 a 154 bytes}. Una reducción global del \textbf{80,2 \%}.
\end{itemize}

[\textbf{PAUSA.}] Para un parque de impresoras que reporta cada 5 segundos durante todo el día, ese 80 \% se traduce directamente en un ahorro del consumo de datos del contribuyente.

\subsection*{Diapositiva 14 — Resultado 4: resiliencia, cero pérdida de datos (15:57 – 17:18)}

[\textbf{SEÑALAR la captura de \textit{Log Capture} y el diagrama de flujo.}]

El cuarto resultado es la \textbf{resiliencia}: el sistema tolera fallas de energía y de red \textbf{sin pérdida de datos}.

Cuando el enlace se pierde, al expirar el \textit{Keep-Alive}, el \textit{broker} publica la "última voluntad" del dispositivo. Como ven en la captura, el \textit{worker} interceptó la alerta \texttt{offline} y registró el evento como "Desconexión Abrupta", \textbf{congelando el último estado válido} sin alterar el histórico. El ciclo completo es: pérdida del enlace, reintentos con \textit{backoff} exponencial, publicación del testamento, registro del evento y restauración del túnel.

Esto se apoya en una propiedad de diseño: el módulo es un \textbf{observador de estado sin memoria tributaria propia}. Toda la información fiscal reside en el periférico y se re-consulta en cada sondeo; el \textit{gateway} no es una fuente de verdad que pueda corromperse.

[\textbf{SEÑALAR la tabla de escenarios.}]

En los \textbf{cuatro escenarios de falla} que evaluamos —corte abrupto de energía, micro-corte Wi-Fi, caída del \textit{broker} y desconexión del periférico— \textbf{la pérdida de datos fue nula}.

\subsection*{Diapositiva 15 — Resultado 5: latencia End-to-End ($T_{E2E}$) (17:18 – 18:38)}

[\textbf{SEÑALAR la tabla de latencias.}]

El quinto resultado es el \textbf{desempeño temporal}. Definimos la latencia de extremo a extremo como la suma de cinco etapas, desde el sondeo físico hasta el render en el navegador. Sobre \textbf{500 iteraciones continuas}, obtuvimos:

[\textbf{SEÑALAR cada fila.}]

\begin{itemize}[label=-]
    \item Sondeo síncrono del S1: \textbf{118 milisegundos} (dentro de la ventana fiscal).
    \item Empaquetado JSON: \textbf{2 milisegundos}.
    \item \textit{Handshake} TLS 1.2: \textbf{45 milisegundos}.
    \item Publicación y tránsito por el \textit{broker}: \textbf{85 milisegundos}, la etapa más variable.
    \item Renderizado del DOM: \textbf{18 milisegundos}.
\end{itemize}

La \textbf{latencia total promedio es de 268 milisegundos, con desviación de 15,3}: apenas el \textbf{5,36 \% de la ventana de sondeo} de 5 segundos. El sistema opera en tiempo real estricto sin acumular retrasos. El análisis de varianza mostró que el \textbf{86,8 \% de la dispersión} proviene de la etapa de red, es decir, del enlace del proveedor de Internet: la incertidumbre no está en el diseño embebido, sino en un factor externo y acotado.

\subsection*{Diapositiva 16 — Resultado 6: OTA de doble objetivo (18:38 – 20:09)}

[\textbf{SEÑALAR el diagrama de flujo de la OTA.}]

El sexto resultado es la \textbf{actualización remota de firmware}, que fue además el aporte más disruptivo del trabajo: un protocolo OTA de \textbf{doble objetivo}.

[\textbf{SEÑALAR la rama izquierda.}]

El primer objetivo es el \textbf{gateway} ESP32. Al recibir la orden, activa una bandera que detiene el sondeo y escribe en la \textbf{partición OTA inactiva} de un esquema dual: si ocurre un corte de energía, el \textit{bootloader} arranca desde la partición intacta.

[\textbf{SEÑALAR la rama derecha.}]

El segundo es el más complejo: \textbf{reescribir in situ el firmware del periférico fiscal}, que no tiene pila de red. El módulo actúa como \textit{host} de programación intermedio: como el binario supera el \textit{heap} libre del ESP32, la descarga no se almacena completa, sino que se procesa en fragmentos de 1024 bytes a través de un \textbf{búfer de trabajo fijo}; cada fragmento viaja con cabecera y redundancia, y se admite hasta 3 reintentos con \textit{ACK}. El cierre \textbf{solo se emite cuando los bytes confirmados igualan la longitud declarada}, así que una transferencia parcial nunca se compromete; y ante cortes de red, la reanudación usa \textbf{cabeceras HTTP Range}, sin reiniciar desde cero.

[\textbf{SEÑALAR la captura de éxito.}] \textbf{Ambos objetivos quedaron validados de extremo a extremo.}

\subsection*{Diapositiva 17 — Resultado 7: rendimiento global (20:09 – 21:31)}

[\textbf{SEÑALAR la tabla de perfilado de memoria.}]

El séptimo resultado es el \textbf{rendimiento global}. Empecemos por la memoria, el recurso más restringido. En el ESP32, la IRAM y el \textit{heap} comparten la \textbf{misma SRAM física}: lo que se libera de una queda disponible para la otra. Tras la reingeniería mediante \texttt{menuconfig} —desactivando la aceleración por IRAM de la pila Wi-Fi—, redujimos la \textbf{IRAM del 78,39 \% al 52,26 \%}: un ahorro de 34.256 bytes, un \textbf{26,13 \% menos}, y ese margen de SRAM fue el que hizo viable la reserva del búfer dinámico que exigen las operaciones de actualización.

[\textbf{SEÑALAR la tabla del subclocking.}]

En el plano energético adoptamos un \textbf{subclocking a 80 MHz} frente a los 240 MHz nominales: como la potencia dinámica es proporcional a la frecuencia, cae a \textbf{un tercio}, y con ella el salto térmico, así que el sistema opera \textbf{más frío}. Lo habilita el \textbf{acelerador AES por hardware}: los 80 MHz sostienen el \textit{handshake} TLS en 45 milisegundos.

[\textbf{SEÑALAR la tabla de SQLite.}]

Y en la nube, la trazabilidad: con \textbf{10.000 registros}, la búsqueda indexada por serial se mantiene \textbf{por debajo de los 12 milisegundos}, y una depuración automática purga lo mayor a 90 días sin bloquear las consultas.

\subsection*{Diapositiva 18 — Resultado 8: comparativa estado previo vs. actual (21:31 – 22:42)}

[\textbf{SEÑALAR la tabla "antes y después". Cierra el bloque de resultados.}]

Para sintetizar los resultados, la comparativa.

[\textbf{SEÑALAR columna izquierda.}]

\textbf{Estado previo:} el técnico conectaba físicamente una PC al puerto serial; observaba \textbf{un equipo a la vez}, sin visión global del parque; el diagnóstico era reactivo, posterior al fallo; no había historial de auditoría; y actualizar firmware exigía intervención en sitio.

[\textbf{SEÑALAR columna derecha.}]

\textbf{Estado actual:} monitoreo remoto y centralizado de \textbf{múltiples equipos en simultáneo}; detección temprana y automática de errores, con alertas; \textbf{bitácora de 90 días} consultable por serial; mantenimiento preventivo con datos reales; y actualización de firmware \textbf{desde la nube}, tanto del \textit{gateway} como del periférico.

[\textbf{PAUSA}] En una frase: de un equipo aislado y observado a mano, a un parque conectado, trazable y administrable, con un 80,2 \% menos de datos y 268 milisegundos de latencia de extremo a extremo.

[\textbf{TRANSICIÓN al Bloque 4}] Con esto concluimos los resultados. Pasamos al impacto y las conclusiones.

\hrulefill

\section*{BLOQUE 4 — Impacto, Conclusiones y Escalabilidad}
\textbf{Tiempo estimado: 4,7 min (acumulado 22:42 – 27:21) $\cdot$ Diapositivas 19 a 22}

\subsection*{Diapositiva 19 — Verificación de cumplimiento de objetivos (22:42 – 23:57)}

[\textbf{SEÑALAR la matriz de objetivos del informe final.}]

Cerramos verificando el cumplimiento de objetivos. El informe incluye una \textbf{matriz de cumplimiento} que contrasta cada objetivo con su estado y su entregable técnico.

Los \textbf{ocho objetivos del plan se alcanzaron}, con un matiz que declaramos con transparencia:

\begin{itemize}[label=-]
    \item \textbf{Arquitectura, firmware, integración y documentación:} logrados.
    \item \textbf{Comunicación MQTT:} logrado. Cliente MQTT v5.0 sobre TLS 1.2, con \textit{bundle} de certificados, \textit{Last Will and Testament} y \textit{User Properties}.
    \item \textbf{Servidor MQTT (recepción, almacenamiento y distribución):} \textbf{parcialmente logrado}. En lugar de desplegar un \textit{broker} propio, se delegó el servicio al \textit{broker} gestionado HiveMQ Cloud; la capa de recepción, almacenamiento y distribución sí se implementó.
    \item \textbf{Dashboard en tiempo real y validación:} logrados. Interfaz reactiva y métricas cuantificadas sobre banco HIL.
\end{itemize}

[\textbf{PAUSA}] Además, \textbf{ampliamos el alcance respecto al plan original} con un objetivo no previsto: el \textbf{módulo OTA de doble objetivo}, que resultó el componente más valioso del trabajo.

\subsection*{Diapositiva 20 — Aporte y beneficio para la organización (23:57 – 24:58)}

[\textbf{SEÑALAR la infografía de beneficios.}]

¿Qué significa esto para The Factory HKA? Cuatro beneficios directos.

Primero, \textbf{visibilidad centralizada del parque}: lo que antes exigía recorrer físicamente cada equipo hoy se observa en un solo panel.

Segundo, \textbf{mantenimiento preventivo basado en datos}: la bitácora de 90 días permite analizar patrones de fallo y anticipar intervenciones, reduciendo tiempos de inactividad en las impresoras fiscales.

Tercero, \textbf{ahorro operativo en red}: el 80,2 \% menos de datos por trama se traduce en menores costos de conectividad a escala de flota.

Y cuarto, \textbf{mantenimiento de firmware sin intervención en sitio}: la actualización remota permite corregir el software tanto del \textit{gateway} como de la impresora desde la nube, algo clave para una empresa con equipos dispersos geográficamente.

\subsection*{Diapositiva 21 — Limitaciones y trabajo futuro (24:58 – 26:17)}

[\textbf{SEÑALAR la lista. El jurado valora la autocrítica.}]

Toda solución tiene límites, y queremos ser claros con los de este trabajo.

\begin{itemize}[label=-]
    \item \textbf{Validación en laboratorio:} se validó sobre el banco HIL. Es un prototipo funcional y documentado, pero \textbf{aún no está desplegado en campo}; el siguiente paso es una prueba piloto con un conjunto reducido de contribuyentes.
    \item \textbf{Escala de la nube:} SQLite es adecuado para un piloto, pero para un parque nacional conviene migrar a un motor con mayor concurrencia, como PostgreSQL.
    \item \textbf{Seguridad y cumplimiento:} la información viaja cifrada, pero antes de un despliegue masivo debe revisarse el cumplimiento normativo sobre custodia de datos fiscales.
    \item \textbf{Dependencia de la red del contribuyente:} el sistema es resiliente, pero el 86,8 \% de la dispersión proviene de ese factor externo.
\end{itemize}

Y como \textbf{trabajo futuro} proponemos: extender la telemetría a los reportes Z; incorporar analítica predictiva sobre el histórico de errores, para anticipar fallos mecánicos; y añadir aprovisionamiento masivo de dispositivos.

\hrulefill

\section*{BLOQUE 5 — Cierre y Agradecimientos}
\textbf{Tiempo estimado: 1,6 min (acumulado 27:21 – 28:57) $\cdot$ Diapositiva 22}

\subsection*{Diapositiva 23 — Cierre, agradecimientos y preguntas (27:21 – 28:57)}

[\textbf{Diapositiva final. Baja el ritmo, cierra con seguridad.}]

[\textbf{PAUSA}] En conjunto, la solución deja las condiciones técnicas para que The Factory HKA centralice la telemetría de un parque nacional y transforme la gestión reactiva del mantenimiento en una práctica basada en datos.

[\textbf{PAUSA. Mira al jurado.}]

Quiero agradecer a la Universidad Simón Bolívar y a The Factory HKA; al profesor Miguel Díaz por su tutoría académica; a los ingenieros Pablo Arguinzones y Luis Muñoz por su acompañamiento industrial; y a todas las personas que me apoyaron durante estos meses.

Con esto damos por concluida la presentación y quedamos atentos a sus preguntas. Muchas gracias.

[\textbf{PAUSA. Espera. No agregues nada más.}]

\hrulefill

\section*{Anexo — Preparación para la ronda de preguntas}

Respuestas de 3 a 5 frases. Se dan en primera persona del plural, salvo si la pregunta es personal.

\paragraph{1. ¿Qué bus síncrono usaron: SPI o I2C?}
SPI. Fue lo que ofrecía el hardware de la impresora fiscal: el periférico expone su interfaz síncrona a través de su circuito integrado interno. No fue una decisión de comparar SPI contra I2C, sino la interfaz que el equipo ya traía; y encajó bien porque el ESP32 integra controlador SPI nativo y el esquema maestro-esclavo es directo.

\paragraph{2. ¿Por qué MQTT y no HTTP o WebSockets puros?}
Porque era el protocolo definido en el alcance de la propuesta y, además, resultó mediblemente más eficiente: redujo el volumen por trama en un 80,2 \% frente a HTTP POST. Su modelo publicación-suscripción desacopla el dispositivo de los consumidores —el dashboard y el \textit{worker} de ingesta reciben la misma telemetría—, y el \textit{Last Will and Testament} permite detectar caídas abruptas. El margen de IRAM no dependió del protocolo: la presión venía de la pila Wi-Fi y de TLS, y se resolvió con \texttt{menuconfig}.

\paragraph{3. ¿Cómo garantizan que no se corrompe la información fiscal?}
Porque el ESP32 solo escucha: no manipula ni escribe la memoria fiscal, únicamente lee el \textit{Status S1} y el \textit{SV2}. No se puede corromper lo que no se modifica. Y en la actualización de firmware, antes de aplicar cualquier cambio se descarga y se almacena el binario de forma segura desde el portal.

\paragraph{4. ¿Qué pasa si se corta la energía a mitad de una OTA?}
En el periférico fiscal —el caso validado— el cierre de la transacción solo se emite cuando los bytes confirmados igualan la longitud declarada: una transferencia parcial nunca se compromete y una nueva orden reinicia el proceso desde cero. En el \textit{gateway}, por diseño, la escritura ocurre sobre la partición OTA inactiva y el \textit{bootloader} arrancaría desde la partición intacta; ese comportamiento está pendiente de confirmar en pruebas.

\paragraph{5. ¿Por qué SQLite y no PostgreSQL?}
SQLite es suficiente para el piloto: con 10.000 registros, la búsqueda indexada se mantiene por debajo de 12 milisegundos. Reconocemos que para un parque nacional conviene PostgreSQL por concurrencia; lo dejamos como trabajo futuro.

\paragraph{6. ¿Por qué ESP32 y no una Raspberry Pi o un PLC?}
Porque es un proyecto IoT, y para este caso el ESP32 es la plataforma indicada: doble núcleo con FreeRTOS, aceleración AES por hardware, bajo consumo y bajo costo, sin sistema operativo de propósito general. Además, es la plataforma que la empresa definió para sus equipos fiscales, lo que da continuidad con su ecosistema. Una Raspberry Pi añadiría un SO completo y mayor superficie de ataque, sin aportar nada que el sistema requiera.

\paragraph{7. ¿Es seguro enviar datos fiscales a la nube?}
Sí. El transporte usa TLS 1.2 de extremo a extremo, que cifra la información en tránsito, con validación de la cadena de confianza mediante \textit{bundle} de certificados. Eso protege los datos frente a interceptación o manipulación.

\paragraph{8. ¿Cómo comprobaron que el problema del bus era por interferencia electromagnética?}
Lo evidenciamos por sus síntomas: sobre el enlace asíncrono aparecían errores de enmarcado y fallos en la verificación LRC, que desaparecieron al migrar a la interfaz síncrona con reloj dedicado. No lo caracterizamos con osciloscopio, y lo reconocemos como mejora: una campaña con instrumentación permitiría cuantificar la inmunidad del enlace.

\paragraph{9. ¿Cómo escala a miles de dispositivos?}
La arquitectura escala por diseño: añadir un dispositivo es añadir una suscripción a un tópico. El \textit{broker} gestiona las sesiones y el \textit{worker} de ingesta se desacopla del servidor web. Los cuellos de botella previsibles serían la base de datos y el aprovisionamiento, atacados con la retención de 90 días.

\paragraph{10. ¿La latencia de 268 ms es aceptable?}
Sí, con holgura: representa el 5,36 \% del intervalo de sondeo de 5 segundos. Además, el 86,8 \% de la dispersión proviene del enlace del proveedor de Internet, no del diseño embebido, lo que significa que en un enlace estable el sistema sería aún más regular.

\paragraph{11. ¿Por qué el objetivo del servidor MQTT quedó parcialmente logrado?}
Porque el plan contemplaba un \textit{broker} propio, y en la arquitectura final se delegó al \textit{broker} gestionado HiveMQ Cloud, que ofrece disponibilidad y seguridad sin administración local. La capa que sí implementamos fue la de recepción, almacenamiento y distribución en el servidor Django. Preferimos declararlo así en lugar de forzar la lectura del objetivo.

\paragraph{12. ¿Cuál fue el aporte más significativo?}
El puente síncrono heterogéneo de reprogramación inalámbrica. Es la capacidad más disruptiva porque traslada el mantenimiento de firmware a la nube y reduce la intervención física en campo, algo crítico para una empresa con equipos dispersos geográficamente.

\paragraph{13. ¿Por qué estaba desactivada la verificación del nombre del host en la OTA, y cómo se resolvió?}
Fue una decisión de compatibilidad del prototipo: en el laboratorio la descarga se servía por HTTP, así que la validación del certificado no se ejercía y la bandera quedó activa. En una revisión de seguridad del código la detectamos y la corregimos: la verificación del nombre del host está activa, la ruta de descarga del periférico valida la cadena de confianza, y el firmware rechaza cualquier URL que no sea HTTPS, de modo que la descarga no puede degradarse a HTTP. Como refuerzo adicional, queda recomendada la firma del firmware con Secure Boot, para que la integridad no dependa únicamente del canal.

\paragraph{14. ¿Qué es HTMX y en qué se diferencia de HTML?}
HTML es el lenguaje de marcado que describe la estructura de la página y, por sí solo, es estático: un enlace o un formulario solo pueden provocar una recarga completa del documento. HTMX es una biblioteca ligera de JavaScript que \textit{extiende} HTML con atributos propios ---\texttt{hx-get}, \texttt{hx-post}, \texttt{hx-target}, \texttt{hx-swap}--- para que cualquier elemento dispare peticiones al servidor y actualice solo el fragmento que cambia, sin recargar la página y sin escribir JavaScript. La diferencia conceptual es qué devuelve el servidor: con HTML puro siempre devuelve documentos completos; con HTMX devuelve fragmentos de HTML que se intercambian en su lugar dentro del DOM. En el \textit{dashboard} lo usamos para la navegación y las vistas parciales ---listas, detalles, formularios---, mientras la telemetría en vivo la lleva el cliente MQTT del navegador sobre el canal WSS; esa combinación da la experiencia de página única sin un \textit{framework} de \textit{frontend} pesado.

\paragraph{15. ¿Por qué validaron sobre emuladores y no sobre la impresora real?}
Por dos razones que se combinaron. La metodológica —como se explica en el informe— es que inducir fallos destructivos sobre un equipo fiscal real arriesga registros tributarios, así que el banco HIL era la vía correcta en cualquier caso. La práctica fue de disponibilidad: el firmware del periférico seguía en desarrollo por parte de la empresa y el equipo solo llegó en la última semana. Con la HKA80 ya disponible, la integración final se demuestra sobre el equipo real en esta presentación.

\paragraph{16. ¿Qué es el DOM y qué significa ``renderizar el DOM''?}
El DOM (\textit{Document Object Model}) es la representación en memoria que el navegador construye de la página: un árbol de objetos donde cada etiqueta —un \texttt{div}, una tabla, un indicador— es un nodo manipulable. La diferencia con el HTML es que el HTML es el texto estático que llega del servidor, mientras el DOM es la estructura viva que el navegador crea a partir de él. Cuando decimos ``renderizado del DOM'' —los 18 ms de la matriz de latencias— nos referimos al tiempo que tarda el navegador en actualizar ese árbol con la telemetría recibida y repintar la interfaz. HTMX y el cliente MQTT.js operan precisamente sobre ese árbol: insertan fragmentos nuevos sin reconstruir la página completa.

\paragraph{17. ¿Cómo transfiere el firmware al periférico si el binario no cabe en la RAM del ESP32?}
Porque el binario nunca se almacena completo: la descarga se procesa en \textit{streaming} a través de un búfer fijo de 1024 bytes que se reutiliza en cada fragmento, así que la huella de memoria es constante e independiente del tamaño del firmware. Cada fragmento viaja por el bus síncrono con cabecera de secuencia y redundancia, con hasta 3 reintentos esperando el \textit{ACK} del periférico. Ante cortes de red la descarga se reanuda con la cabecera HTTP \texttt{Range} desde el último byte confirmado, y el cierre solo se emite cuando los bytes recibidos igualan la longitud declarada por el servidor.

\paragraph{18. ¿Por qué el firmware no se descarga por MQTT, si el sistema está basado en MQTT?}
Porque MQTT es un protocolo de mensajería, no de transferencia de archivos: un binario de $\sim$900 KB habría obligado a fragmentarlo en cientos de mensajes y a implementar a mano la secuenciación, el reensamblado y la reanudación —reinventar sobre MQTT lo que HTTP ya ofrece con \texttt{Range}, \texttt{Content-Length} y \textit{streaming}—. Además, el binario atravesaría el \textit{broker}, que limita y factura el tráfico, en lugar de descargar directo del servidor. La división que usamos es el patrón estándar en IoT: MQTT para el plano de control —la orden OTA sí llega por MQTT— y HTTPS para el plano de datos masivo. Es la misma lógica de eficiencia del informe: MQTT gana en mensajes pequeños y frecuentes, HTTP gana en transferencias grandes y ocasionales.

\paragraph{19. ¿Qué es la IRAM y por qué liberarla habilitó la OTA?}
La IRAM es la porción de SRAM donde reside el \textit{código} que debe ejecutarse mientras el caché de flash está deshabilitado —interrupciones, escrituras a flash— o que corre por velocidad; la aceleración por IRAM de la pila Wi-Fi copiaba funciones de red ahí y consumía el 78 \%. En el ESP32 la IRAM y el \textit{heap} comparten la misma SRAM física, así que al desactivar esa opción en \texttt{menuconfig}, el código de Wi-Fi volvió a flash y los 34 KB liberados quedaron disponibles como memoria de trabajo: el margen que necesitaban el búfer de descarga y el \textit{handshake} TLS de la OTA.

\paragraph{20. ¿Qué es AES y qué papel juega en el sistema?}
AES (\textit{Advanced Encryption Standard}) es el algoritmo de cifrado simétrico estándar —la misma clave cifra y descifra— sobre bloques de 128 bits. En nuestro sistema es el algoritmo que cifra el canal TLS: el \textit{handshake} negocia una clave de sesión con criptografía asimétrica y, desde ahí, AES cifra el volumen de datos porque es mucho más rápido. Además, el ESP32 incluye un acelerador AES por hardware —un bloque de silicio dedicado, no el núcleo—, y por eso el \textit{handshake} se sostiene en 45 ms incluso con el CPU a 80 MHz.

\subsection*{Preguntas proyectando el perfil del tutor académico (comunicaciones inalámbricas, antenas, electrónica, IoT)}

\paragraph{21. ¿Por qué Wi-Fi y no un enlace celular —NB-IoT, LTE-M— o LoRa para el WAN?}
Wi-Fi fue la opción coherente con el despliegue: los puntos de venta ya tienen infraestructura de red local, el costo recurrente es cero —no hay SIM ni plan de datos por equipo— y el módulo de transmisión que la empresa fabrica ya incorpora el ESP32 con radio Wi-Fi. Reconocemos que NB-IoT daría cobertura profunda donde no haya LAN y LoRa daría menor consumo; la ruta de trabajo futuro es un respaldo celular para sitios sin Wi-Fi, y la arquitectura no cambia porque MQTT y TLS corren igual sobre TCP/IP.

\paragraph{22. ¿Caracterizaron el enlace de radio? ¿RSSI, presupuesto de enlace?}
El RSSI viaja como telemetría: cada publicación MQTT lo inyecta en las \textit{User Properties}, así que el \textit{dashboard} ve la calidad de señal de cada equipo —es una métrica operativa, no solo de laboratorio. Lo que no hicimos fue una campaña formal de presupuesto de enlace ni mapa de cobertura por ubicación; lo registramos como mejora, y justamente el RSSI ya en la plataforma es lo que permitiría hacerla con datos de producción.

\paragraph{23. ¿Cómo atraviesa la señal Wi-Fi el gabinete metálico de la impresora?}
Punto válido: el prototipo del banco HIL opera con el ESP32 en tarjeta de desarrollo, no dentro del chasis de la HKA80, así que el efecto jaula de Faraday no se ejercitó en las pruebas. Para el producto hay dos rutas conocidas: ubicar la antena en una ventana no metálica del gabinete o migrar a un módulo con conector de antena externa. El RSSI que ya telemetramos sería la métrica de aceptación de esa integración mecánica.

\paragraph{24. ¿Cómo se comporta el sistema con el espectro de 2,4 GHz congestionado, como en un centro comercial?}
La telemetría es de muy bajo \textit{duty cycle} —una trama de $\sim$150 bytes cada 5 segundos—, así que la contienda de canal degrada la latencia antes que la integridad: ese \textit{jitter} es justamente lo que medimos como el 86,8 \% de dispersión en la etapa de red. QoS 1 retransmite ante pérdida y el firmware reintenta con retardo de cortesía para no contribuir a la congestión. No hicimos un \textit{site survey} de canales; las fallas se validaron inyectándolas directamente.

\paragraph{25. ¿Cómo saben que el problema del bus asíncrono era EMI y no un simple desajuste de baudios?}
Es una distinción justa: los síntomas —errores de enmarcado y fallos de LRC intermitentes— son compatibles con ambas causas, y nuestra evidencia fue por el lado del síntoma, sin osciloscopio. Lo importante es que la migración al bus síncrono elimina la categoría completa del problema: al compartir reloj ya no hay deriva posible, y el enlace estructurado reduce la susceptibilidad al ruido —así que la decisión era correcta bajo cualquiera de las dos hipótesis.

\paragraph{26. ¿Qué niveles eléctricos maneja el bus síncrono y cómo protegen el acoplamiento con el equipo fiscal?}
El bus trabaja a los niveles lógicos que expone el circuito interno del periférico, con el ESP32 como maestro generando reloj y habilitación. En el prototipo el acoplamiento es directo porque es la interfaz que el hardware ofrece; para producción la recomendación clásica es el acondicionamiento de líneas —e incluso optoacopladores para aislar masas—, que es donde entra el diseño de hardware de los módulos de la empresa.

\paragraph{27. ¿Por qué un \textit{checksum} XOR y no un CRC?}
Porque la redundancia no la elegimos nosotros: la hereda el protocolo fiscal. El periférico valida con LRC —el XOR longitudinal es el estándar de estos equipos— y nuestro empaquetado OTA replica el mismo esquema para que el esclavo no implemente dos verificadores. Conocemos sus límites: un XOR no detecta un número par de errores en la misma posición de bit ni reordenamientos; un CRC-8 sería el \textit{upgrade} natural, pero exigiría un cambio de firmware en el periférico fiscal.

\paragraph{28. El bus corre a 100 kHz, pero el sondeo del S1 toma 118 ms. ¿Por qué no unos 13 ms teóricos?}
Porque los 118 ms no son solo tiempo de bits: incluyen el ciclo completo de la transacción —la encuesta ENQ, el tiempo de procesamiento del periférico para armar la trama, los márgenes entre fases y el desempaquetado—. El reloj gobierna la transferencia de datos, pero el tiempo total lo domina el protocolo de la impresora, no el bus; aun así quedó dentro de la ventana fiscal de 200 ms.

\paragraph{29. ¿De dónde se alimenta el módulo y cómo manejan los picos de corriente del radio?}
El módulo se alimenta de la fuente interna del equipo; el punto delicado de un ESP32 es la ráfaga de corriente en transmisión Wi-Fi, que puede superar los 300-400 mA. Lo que mitiga nuestra configuración es el \textit{subclocking} a 80 MHz: la demanda instantánea se vuelve más uniforme y las transitorias de corriente se reducen —ese fue uno de los criterios para adoptar la frecuencia baja. La verificación instrumental de la forma de onda de consumo quedó como mejora.

\paragraph{30. El \textit{subclocking} reduce la potencia dinámica del núcleo, pero el PA de RF es fijo. ¿Cuánto baja el consumo total real?}
Correcto: el factor 3 de reducción aplica a la potencia dinámica del núcleo —proporcional a $f$— y el radio opera a potencia constante. Por eso el informe lo enmarca como decisión térmica: el salto térmico unión-ambiente del SoC cae a un tercio, que era lo que calentaba el gabinete en operación continua. El consumo total del módulo se reduce menos, porque la fracción RF no depende del \textit{clock}; el objetivo era estabilidad térmica 24/7, no ahorro energético.

\paragraph{31. Con QoS 1 pueden llegar mensajes duplicados. ¿Cómo deduplican?}
Por diseño no hace falta deduplicar paquetes: el \textit{worker} no persiste mensajes, persiste eventos —compara cada trama con el último estado guardado y solo registra la diferencia. Un duplicado del mismo S1 produce una diferencia vacía, así que la bitácora queda idempotente de forma natural. Es telemetría de estado, no de flujo de eventos.

\paragraph{32. ¿Cómo estampan la hora real de cada evento?}
El \textit{timestamp} que viaja en el JSON es relativo al arranque del ESP32 —sirve para secuenciar—, y la estampación de hora real la hace el \textit{worker} de Django al recibir y persistir. Es una limitación conocida: para una bitácora fiscal ideal los eventos deberían estamparse en origen con NTP, que el ESP32 soporta nativamente; queda como mejora, porque el diseño actual confía la cronología al instante de ingesta.

\paragraph{33. ¿En cuánto tiempo detectan que un equipo se cayó?}
La detección de caída abrupta la da el \textit{Last Will and Testament}: el \textit{broker} publica la alerta cuando la sesión expira, así que el tiempo de detección lo gobierna el \textit{keep-alive} de la sesión. Hay un segundo nivel observable: la ausencia del reporte de 5 segundos se ve en el \textit{dashboard} como dispositivo sin telemetría reciente; una alerta proactiva de ``sin datos en N segundos'' es la mejora natural.

\paragraph{34. El \textit{broker} está alojado en Europa. ¿No impacta eso la latencia?}
Sí, y es parte de por qué la etapa de red concentra el 86,8 \% de la dispersión: el RTT transatlántico está dentro de los 85 ms de publicación y tránsito. El criterio se cumple igual —268 ms es el 5,36 \% de la ventana de 5 s—, pero para producción un \textit{broker} regional o un plan con región cercana reduciría esa etapa directamente.

\paragraph{35. ¿Por qué sondear cada 5 segundos en lugar de reporte por evento?}
Porque el periférico no puede empujar: el protocolo fiscal es maestro-esclavo de solicitud-respuesta, la impresora solo contesta cuando se le encuesta. El intervalo de 5 s es un compromiso: resolución suficiente para capturar eventos —un corte Z, un error, una desconexión— sin saturar el canal, y como el ciclo completo toma 268 ms, el bus queda libre el 95 \% del tiempo para la operación normal del equipo.

\paragraph{36. ¿La sesión MQTT es persistente? ¿Se pierde telemetría durante una caída?}
La telemetría no es un flujo que se acumula: cada mensaje es una fotografía del estado completo, así que lo que interesa es el último valor, no el historial en tránsito. Al reconectar, el equipo publica el estado vigente y el \textit{worker} reconstruye la continuidad; el LWT marca el intervalo de caída. Encolar mensajes viejos con sesión persistente añadiría complejidad para recuperar datos ya desactualizados.

\hrulefill

\section*{Consejos de exposición}

\begin{itemize}[label=-]
    \item \textbf{Ensaya en voz alta con cronómetro.} Si te pasas de 32 minutos, recorta la diapositiva 9 (justificación tecnológica) y la 17 (rendimiento global), que son las más prescindibles.
    \item \textbf{No leas las diapositivas.} Ellas llevan cifras y diagramas; tu voz cuenta la historia que las conecta.
    \item \textbf{Ancla las cifras fuertes} (80,2 \%; 78,39 \% $\rightarrow$ 52,26 \%; 268 ms; cero pérdida de datos) diciéndolas despacio y con una pausa después.
    \item \textbf{Mantén el plural de forma consistente} en la parte técnica. Si se te escapa un "yo", sigue adelante sin corregirte en voz alta.
    \item \textbf{Cuando señales un diagrama, recórrelo con el puntero} y nombra cada nivel antes de explicar su función.
    \item \textbf{Reserva 30 segundos de margen} antes de la primera pregunta.
    \item \textbf{Ten a mano una copia impresa} del informe con la matriz de objetivos y las tablas de resultados.
\end{itemize}

\end{document}